\section{Leer los resultados y convertirlos en una recomendación}\label{sec:lectura-final}

\subsection{Antes de sumar euros: qué cuenta cada concepto}
El catálogo separa beneficio bruto de la marca, VM y CVP. En la simulación,
el VM depende de la edad del coche vendido:
\begin{align}
 \mathrm{CVP}&=\mathrm{beneficio\ bruto}-\mathrm{VM},\\
 R&=\mathrm{CVP}-C_{\mathrm{stock}}-C_{\mathrm{transfer}}
       -C_{\mathrm{CO_2}}-C_{\mathrm{arancel}}-C_{\mathrm{cambio}}.
\end{align}
$R$ es contribución después de los costes modelizados, no beneficio neto contable
de una empresa con todos sus gastos generales. Un coste de cambio negativo
representa un ahorro y, al restarlo, aumenta el resultado.

No se vuelve a restar VM después de calcular CVP. Tampoco se descuentan otra vez
las ventas perdidas: sus ingresos ya no forman parte de las ventas realizadas.
El coste de oportunidad se informa aparte. Los aranceles y el cambio adicional
se imputan a las ventas de cohortes de nueva importación; el coste incremental
que permanece en el stock final se registra, pero no se convierte en ingreso.

\tabla{unidad.tex}{Economía de una unidad joven por segmento, antes de stock y
otros recargos. Todos los importes son euros por coche; el coste de transfer
es el 20\,\% del beneficio bruto, no del CVP.}
{tab:unidad}{\archivo{resultados/escenarios/tablas/economia_unitaria_segmento.csv}}

Un MSUV joven deja \num{5887} euros de CVP antes de costes logísticos; mantener
un coche nuevo 15 días cuesta 30 euros. No es una comparación que pueda
extenderse sin límite: el gasto de stock de la marca se detiene en 360 euros y
el VM de una unidad vieja puede ser diferente.

En B el CVP joven es sólo 125 euros antes de stock, frente a una prima CO$_2$
de \num{2555}. En C, el CVP joven de \num{1367} pasa a $-585$ euros al cargar
CO$_2$, antes de arancel y stock. Una venta adicional no siempre mejora la
contribución. Por eso, una vez confirmados los costes del escenario, se frenan
\emph{nuevas reposiciones} marginalmente negativas. No se deja de vender
automáticamente el inventario ya comprometido.

\subsection{Resultado agregado de la política recomendada}
La tabla~\ref{tab:resultado} utiliza medias de 1.000 réplicas. En particular,
\emph{pedidos A013} cuenta decisiones de ese año, mientras \emph{llegadas}
incluye coches comprometidos en A012. No tienen por qué coincidir.

\tabla{resultados.tex}{Política recomendada: resultados medios A013, total de la
marca. Las coberturas usan demanda mensual media de cada réplica.}
{tab:resultado}{\archivo{resultados/escenarios/tablas/tabla5_plan_recomendado_a013.csv}}

\paso{Central.}
Se simulan unas \num{18539} ventas, \num{478} compras perdidas y un servicio
medio del 97,6\,\%. La contribución media es 62,88 M\euro{}. El stock final
medio es \num{2221} coches. La demanda media simulada de \num{19017} no tiene
por qué igualar exactamente la previsión puntual de \num{19077}: procede de
una muestra finita de trayectorias y del redondeo a unidades.

\paso{Pesimista.}
La contribución cae a 3,09 M\euro{} y el servicio al 77,8\,\%. No se debe sólo
a una demanda menor. CO$_2$ y el recargo de importación hacen negativas nuevas
reposiciones de B, C, SSUV y S bajo la regla marginal, y se incorpora la
interrupción de suministro. El deterioro de servicio es una consecuencia
importante de la política económica, no una cifra que pueda ocultarse con el
cumplimiento del central.

\paso{Optimista.}
Se simulan unas \num{24952} ventas y 91,37 M\euro{} de contribución media.
Sin embargo, se pierden unas \num{3493} compras: la demanda mejora más deprisa
que la capacidad de modificar llegadas. El ahorro de divisa, registrado como
coste negativo, explica que la contribución final pueda superar el CVP antes
de esos ajustes. No se han restado los costes dos veces ni creado ventas extra.

\tabla{riesgo.tex}{Dispersión del resultado total y del servicio anual.
P10--P90 describe variación dentro del escenario, no probabilidades de escenarios.}
{tab:riesgo}{\archivo{resultados/escenarios/tablas/tabla5_plan_recomendado_a013.csv}}

El central tiene P10--P90 de contribución de 56,45--69,47 M\euro{}. Es un rango
condicionado a los supuestos y a la distribución de errores elegida. El error
estándar de su media de Monte Carlo es aproximadamente 0,159 M\euro{}, una
medida de precisión de simulación que no equivale a ese rango de resultados.
No observar pérdidas en las 1.000 réplicas finales no demuestra que su
probabilidad real sea cero.

\grafica{plan_recomendado_ventas_stock_a013.png}
{Ventas y stock del plan recomendado: mediana y P10--P90 por escenario.
Las bandas reflejan variación mensual de una misma política inicial y sus
revisiones. El stock representado no incluye coches en camino.}
{fig:resultado-mensual}{\archivo{escenarios.ipynb}; 1.000 réplicas por escenario}

\subsection{No todos los segmentos desempeñan el mismo papel}
\tabla{segmentos.tex}{Resultado central por segmento: medias de demanda, ventas,
contribución, stock final y servicio. Las sumas pueden diferir por redondeo.}
{tab:segmentos}{\archivo{resultados/escenarios/tablas/tabla5_plan_recomendado_a013.csv}}

MSUV aporta aproximadamente 31,91 M\euro{}, cerca de la mitad del resultado
central. SSUV y CSUV aportan 14,86 y 11,77 M\euro{}, respectivamente. La
disponibilidad de esos segmentos tiene más impacto absoluto que una mejora
pequeña en un segmento de bajo volumen.

B apenas aporta 0,02 M\euro{} en el central: mantener su objetivo de servicio
no convierte su economía en robusta. C conserva un inventario inicial elevado;
una cobertura alta no justifica encargar más coches. S combina bajo volumen,
alta dispersión relativa y stock medio elevado frente a la demanda. Su cobertura
debe leerse como señal de exposición, no como garantía de calidad de previsión.

\subsection{Aislar el CO\texorpdfstring{$_2$}{2} sin confundir escenarios}
Para medir únicamente la prima, se toman las mismas ventas y flujos centrales
y se resta la prima del catálogo a cada venta. No se reoptimiza la política
en esta comparación.

\tabla{co2.tex}{Prima CO$_2$ aislada sobre los mismos flujos centrales.
Importes en M\euro{}; efecto negativo significa reducción del resultado.}
{tab:co2}{\archivo{resultados/escenarios/tablas/efecto_co2_aislado_a013.csv}}

La reducción media aislada es 55,01 M\euro{}, dejando 7,87 M\euro{} sobre esos
mismos flujos. El pesimista completo da 3,09 M\euro{} porque también cambia
el volumen servido, los pedidos, los aranceles y el suministro. No sería
correcto atribuir toda la diferencia central--pesimista a CO$_2$.

\subsection{Qué significa equivocarse por arriba o por abajo}
\begin{description}[style=nextline]
\item[Pasarse.] El exceso inmoviliza stock, consume días de coste y puede exigir
más VM al venderse viejo. La comparación debe considerar los costes posteriores
del inventario final, no sólo diciembre.
\item[Quedarse corto.] Se pierde la contribución de una venta potencial o se
incurre en el coste de conseguir un transfer. El margen no obtenido debe
descontar también CO$_2$, arancel y otros costes que no se pagan al no vender.
Puede resultar negativo en una venta marginal inviable.
\end{description}
El simulador no calcula un presupuesto de publicidad óptimo ni una respuesta de
demanda a cada euro de descuento. «Empujar» un segmento significa aquí priorizar
abastecimiento o revisar su política comercial, no asegurar que cualquier campaña
adicional sea rentable.

\subsection{La sensibilidad: comprobar qué hipótesis sostiene la decisión}
Las pruebas cambian un supuesto cada vez y utilizan las mismas 200 réplicas del
piloto. Por tanto, no hay que comparar mecánicamente sus importes con la media
final de 1.000 réplicas. La referencia central de la política final con esas
200 réplicas es 63,41 M\euro{}, frente a 62,88 M\euro{} con 1.000.

\begin{table}[htbp]\centering
\caption{Ejemplos de sensibilidad, 200 réplicas comunes.
Los importes son medias condicionales A013 y no se comparan con otra muestra.}
\label{tab:sensibilidad}
\begin{tabularx}{\linewidth}{@{}Yrr@{}}\toprule
Prueba & \makecell{Resultado\\(M\euro)} & \makecell{Servicio\\medio (\%)} \\\midrule
Central: referencia de las mismas 200 réplicas & 63,41 & 97,7 \\
Central: transfer del 25\,\% del déficit & 63,63 & 98,3 \\
Central: transfer suficiente, cota favorable & 64,29 & 100,0 \\
Pesimista: sin arancel adicional & 7,78 & 78,1 \\
Pesimista: arancel de 1.000 \euro/coche & $-1{,}50$ & 78,0 \\
Pesimista: no frenar reposición negativa & $-0{,}55$ & 98,6 \\
Central: demanda CSUV un 25\,\% menor & 60,34 & 97,9 \\
Central: demanda CSUV un 25\,\% mayor & 66,32 & 97,3 \\\bottomrule
\end{tabularx}
\fuente{\archivo{resultados/escenarios/tablas/sensibilidades_plan_recomendado.csv};
referencia central a partir del piloto y los $z$ seleccionados}
\end{table}

El ejemplo del arancel muestra una incertidumbre de supuestos que no aparece
en la banda aleatoria del escenario base: con otro recargo, el pesimista puede
tener contribución negativa. La alternativa sin freno económico mejora el
servicio, pero pierde dinero. Es un conflicto de objetivos que debe comunicarse.

Las pruebas también cubren llegadas iniciales, edades, VM de CSUV viejo,
correlaciones, efecto de producto, dispersión de CSUV, sesgo residual y nivel
de continuidad A014. No se les asignan probabilidades.

\subsection{Por qué se prolonga el análisis más allá de diciembre}
Los pedidos de octubre--diciembre de A013 llegan en enero--marzo de A014.
No seguirlos puede premiar una política que mejora artificialmente el cierre y
desabastece el trimestre siguiente.

La prolongación conserva el nivel desestacionalizado del último trimestre y el
perfil mensual de mercado A009--A010. No repite el adelanto o el lanzamiento en
cada enero. Al no conocerse el calendario ni la economía A014, se presenta como
sensibilidad de continuidad y se compara también un nivel un 10\,\% menor o mayor.

\tabla{continuidad.tex}{Comparación del plan base y recomendado al extender
hasta junio de A014. Resultado medio tras costes modelizados.}
{tab:continuidad}{\archivo{resultados/escenarios/tablas/comparacion_continuidad_politicas.csv}}

La mejora frente al plan base se mantiene bajo esa continuidad: el central
recomendado acumula 93,00 M\euro{} en 18 meses frente a 92,66 M\euro{} del base.
No se añade un precio de venta ficticio por todo el stock que queda ni se
demuestra con ello un óptimo de horizonte infinito.

\subsection{La recomendación operativa, expresada correctamente}
\begin{enumerate}
\item \textbf{Partir del plan central y los colchones seleccionados.} El plan
inicial asigna \num{16144} coches a pedidos decididos en A013 y prevé
\num{18114} llegadas en A013; son flujos distintos. Los anexos muestran mes de
decisión y mes de recepción.
\item \textbf{Revisar mensualmente sólo lo modificable.} Los pedidos realmente
ejecutados después de observar ventas no tienen por qué sumar lo mismo que el
plan inicial. En el central su media es de unos \num{15756} coches.
\item \textbf{Vigilar segmentos, no sólo el total.} Priorizar la disponibilidad
de los segmentos de mayor contribución; vigilar especialmente la incertidumbre
de CSUV, el stock de C y la fragilidad económica de B.
\item \textbf{Separar señales comerciales y restricciones físicas.} Un descenso
de ventas con la red agotada no demuestra menor demanda. Un nuevo coste confirmado
puede justificar frenar pedidos aunque exista demanda.
\item \textbf{Revisar las hipótesis que dominan el resultado.} Disponibilidad de
transfers, aranceles, prima CO$_2$, fechas de producto y tamaño del adelanto no
deben quedar escondidos en una cifra final.
\end{enumerate}

\tabla{revision.tex}{Umbrales orientativos del cociente de sumas de tres meses:
ventas observadas / previsión. Son referencias de seguimiento, no garantías.}
{tab:revision}{\archivo{resultados/escenarios/tablas/umbrales_revision_a013.csv}}

Los percentiles se calculan sobre el mismo cociente utilizado en la regla, sin
unir años de validación distintos. CSUV utiliza un proxy de SSUV/MSUV ampliado,
no un umbral propio validado. El rango muy amplio de S y el rango inferior a uno
de MD también recuerdan los errores históricos y la muestra limitada: las señales
requieren diagnóstico, no una orden automática de aumentar o reducir producción.

\begin{idea}{Respuesta global a las cinco preguntas del caso}
Se observa crecimiento de clientes en el total de A012, distinto de la evolución
de la gama comparable. Se obtiene una previsión A013 evaluada retrospectivamente,
con una limitación específica para CSUV. Los escenarios muestran rango de ventas,
stock y contribución. El plan separa pedidos de llegadas y revisa a tres meses.
La política recomendada cumple el servicio medio central acordado y mejora la
contribución frente al base, con supuestos y riesgos explícitos.
\end{idea}
\FloatBarrier
